\documentclass[11pt,a4paper]{article}
\usepackage[utf8]{inputenc}
\usepackage{amsmath,amsfonts,amssymb}
\usepackage{geometry}
\geometry{margin=1in}
\usepackage{hyperref}
\usepackage[many]{tcolorbox}
\usepackage{enumitem}
\usepackage{fancyhdr}

\pagestyle{fancy}
\fancyhf{}
\fancyfoot[L]{Universitas Bengkulu - Teknik Elektro}
\fancyfoot[R]{\thepage}
\def\footrule{{\color{gray}\hrule}}

\title{\vspace{-2cm}\textbf{Case Method 2\\Persamaan Diferensial - Minggu 5: PDB Orde 2 Non-Homogen \& Aplikasi RLC}}
\author{Novalio Daratha \& Muhammad Arfan\\Program Studi Teknik Elektro\\Universitas Bengkulu}
\date{2026}

\begin{document}

\maketitle

\begin{tcolorbox}[colback=red!5!white,colframe=red!75!black,title=Instruksi Tugas]
Kumpulkan jawaban Case Method ini pada awal perkuliahan minggu berikutnya. Tunjukkan semua langkah pemikiran, analisis rangkaian, dan penyelesaian matematis Anda secara detil dan sistematis.
\end{tcolorbox}

\section*{Daftar Soal (Level Kognitif C1 - C6)}

\begin{enumerate}[label=\textbf{\arabic*.}]
    \item \textbf{(C1 - Mengingat)} Sebutkan bentuk tebakan solusi partikular $y_p(x)$ menggunakan Metode Koefisien Tak Tentu apabila fungsi $f(x) = 5e^{3x} + 2x^2$.
    \item \textbf{(C2 - Memahami)} Jelaskan mengapa solusi umum persamaan diferensial non-homogen merupakan penjumlahan dari solusi homogen dan solusi partikular. Apa arti fisis dari masing-masing komponen solusi tersebut dalam konteks sistem dinamik (misalnya rangkaian RLC)?
    \item \textbf{(C3 - Mengaplikasikan)} Tentukan solusi partikular dari persamaan diferensial berikut menggunakan metode koefisien tak tentu: $y'' + 3y' + 2y = 4e^{-3x}$.
    \item \textbf{(C4 - Menganalisis)} Diberikan persamaan diferensial $y'' - y' - 2y = 3e^{-x}$. Analisislah mengapa tebakan awal $y_p = A e^{-x}$ tidak akan berhasil, dan tentukan bentuk tebakan yang benar serta solusi partikularnya!
    \item \textbf{(C5 - Mengevaluasi)} Sebuah rangkaian seri RLC dihubungkan dengan sumber tegangan AC $E(t) = 100 \sin(377t)$. Diketahui $R = 10 \, \Omega$, $L = 0.1$ H, dan $C = 10^{-3}$ F. Evaluasi respons tunak (steady-state response) dari arus pada rangkaian tersebut menggunakan PDB non-homogen! Apakah resonansi dapat terjadi?
    \item \textbf{(C6 - Mencipta)} Rancanglah sebuah sistem elektromekanis atau rangkaian kelistrikan yang dapat dimodelkan oleh persamaan diferensial non-homogen $x'' + 2x' + 5x = 10\cos(2t)$. Tentukan parameter fisis sistem Anda dan temukan respons total sistem jika pada saat awal sistem berada dalam keadaan diam ($x(0)=0, x'(0)=0$).
\end{enumerate}

\vfill
\begin{center}
\textbf{Semoga berhasil!}
\end{center}

\end{document}
